\section{Introduction}

\textbf{Research Background:} Unsupervised learning is a pivotal area in artificial intelligence, focusing on identifying significant patterns within large datasets lacking labels. A central challenge in this domain is representation learning \cite{bengio2013representation, lecun2015deep}, which involves deriving meaningful features from raw data inputs. Notable methods such as generative adversarial networks (GANs) \cite{goodfellow2014generative} and vector quantized variational autoencoders (VQ-VAE) \cite{van2017neural} have achieved substantial success in tasks like high-fidelity image generation and domain adaptation. However, these models often exhibit limitations in maintaining consistency across different data distributions and commonly necessitate extensive supervision with labeled data.

Achieving disentangled representations is a critical goal in representation learning, where distinct data characteristics are isolated in separate dimensions. These representations are essential for tasks such as image generation and transfer learning, where understanding data structures without specific supervised tasks is crucial. Despite advances, current VQ-VAE implementations encounter difficulties in handling multi-modal data and capitalizing on unlabeled information to enhance generalization and flexibility.

\textbf{Research Motivation:} Although current methodologies demonstrate considerable strength, they often fail to maintain reliable performance across varied modalities without significant supervision, limiting their application in diverse real-world contexts. There exists a gap in methods that can independently fine-tune their parameters to adapt across different modalities without relying heavily on labeled data. Incorporating self-supervised learning presents a promising opportunity to improve the adaptability and robustness of models, thereby addressing these limitations in unsupervised learning.

\textbf{Methodology Overview:} To address these challenges, we introduce a novel Multi-Modal Vector Quantized Variational AutoEncoder (VQ-VAE) architecture enhanced with self-supervised learning (SSL) features. This approach includes a harmonizer module that facilitates the alignment and transformation of data representations across multiple modalities. Utilizing SSL, the model dynamically refines its internal representations, thereby enhancing image reconstruction quality and adaptability to new domains. This comprehensive strategy allows the model to learn efficiently from unlabeled data, improving generalization without extensive supervision.

\textbf{Contributions:} Our principal contributions are as follows:
\begin{itemize}
    \item We develop a novel Multi-Modal VQ-VAE architecture, incorporating a harmonizer module that aligns representations across diverse data modalities, thereby advancing multi-modal learning.
    \item Our approach integrates self-supervised learning techniques to iteratively refine model parameters, optimizing the embedding space with minimal dependence on labeled data.
    \item We exhibit significant enhancements in image reconstruction fidelity and domain adaptability, validated through extensive empirical evaluations on CIFAR-10.
    \item Our work advances unsupervised learning by integrating modality-specific feature extraction and embedding generalization into a unified framework.
\end{itemize}
